\documentclass[11pt,openany]{book}
% Trade paperback trim size. Change paperwidth and paperheight for other formats.
\usepackage[paperwidth=5.5in,paperheight=8.5in,inner=0.85in,outer=0.65in,top=0.8in,bottom=0.9in]{geometry}
\usepackage{lmodern}
\usepackage{microtype}
\usepackage{titlesec}
\usepackage{fancyhdr}
\usepackage{setspace}
\usepackage{emptypage}

\setstretch{1.12}
\setlength{\parindent}{1.4em}

% Chapter openings: small caps label, big number, then the chapter name, set low on the page.
\titleformat{\chapter}[display]
  {\normalfont\centering}
  {\small\scshape\MakeLowercase{Chapter}\ \thechapter}
  {1ex}
  {\LARGE\itshape}
  [\vspace{0.5ex}\rule{2cm}{0.4pt}]
\titlespacing*{\chapter}{0pt}{2.2in}{0.6in}

\pagestyle{fancy}
\fancyhf{}
\fancyhead[CE]{\small\scshape Author Name}
\fancyhead[CO]{\small\scshape The Salt Road}
\fancyfoot[C]{\small\thepage}
\renewcommand{\headrulewidth}{0pt}
\fancypagestyle{plain}{\fancyhf{}\fancyfoot[C]{\small\thepage}}

% \scenebreak marks a jump in time or place inside a chapter.
\newcommand{\scenebreak}{\par\medskip\begin{center}*\quad*\quad*\end{center}\medskip\noindent}

% The first words of each chapter are set in small caps.
\newcommand{\opening}[1]{\noindent\textsc{#1}}

\begin{document}

\frontmatter
\thispagestyle{empty}
\begin{center}
\vspace*{1.8in}
{\Huge\scshape The Salt Road\par}
\vspace{0.3in}
{\large\itshape A Novel\par}
\vfill
{\large Author Name\par}
\vspace{0.5in}
\end{center}

\clearpage
\thispagestyle{empty}
\vspace*{\fill}
{\small\noindent Copyright \copyright\ 2026 Author Name. All rights reserved.\par
\medskip\noindent This is a work of fiction. Names, characters and places are products of the author's imagination.\par
\medskip\noindent ISBN 000-0-00-000000-0\par}

\clearpage
\thispagestyle{empty}
\vspace*{2in}
\begin{center}\itshape For the people who waited.\end{center}

\mainmatter

\chapter{The Last Ferry}
\opening{The ferry left without her} at a quarter past six, and Ines watched its lights shrink until the fog took them. Replace this paragraph with the opening of your story. Keep the first line strong: it is the one every reader sees.

Paragraphs after the first are indented automatically. Dialogue works as usual.
``You'll want to wait inside,'' the ticket clerk said. ``Nothing else is coming tonight.''

\scenebreak
By morning the harbour had changed. Use a scene break like the one above whenever the point of view, place or time shifts.

\chapter{Inland}
\opening{Three days east of the coast} the road turned to gravel. Continue writing your chapters in the same pattern, one \verb|\chapter| per chapter.

\chapter{What the River Kept}
\opening{Nobody in the village} would say who had built the bridge. Add as many chapters as you need.

\backmatter
\chapter*{Acknowledgements}
Thank your editor, early readers and anyone who helped the book happen.

\chapter*{About the Author}
Author Name lives by the sea. Write two or three sentences in the third person.

\end{document}
